\subsection{The terminal pair at a minimum-two odd root}
\label{proof:Q5-C192}
Use the original finite-space model and messages of
Section~\ref{sec:odd-common-calculus}. The residual lists are
\[
 (2,a,b,s_1,t_1,\ldots,s_k,t_k),\qquad
 3\le a<b<s_1<t_1<\cdots<s_k<t_k,\qquad (2,t),\quad t\ge3.
\]
All lists are finite; the inner and outer lengths \(L,M\ge1\) are
independent. Take exactly the minimum-two directions of
Theorem~\ref{thm:odd-two:roots}: for \(L=1\), the root factor is
\(f_2h_af_b\prod h_{s_i}f_{t_i}\) and the terminal is
\(W_2\bar f_t\). For \(L\ge2\), the terminal is \(\bar f_2W_t\);
the root factor is \(W_2\bar f_3W_b\prod W_{s_i}\bar f_{t_i}\)
at \(a=3\), and \(W_2W_a\bar f_b\prod W_{s_i}\bar f_{t_i}\)
at \(a\ge4\). Every inner edge points terminal-to-root and every
outer edge root-to-receiver. Thus these directions are fixed before
the original host is chosen.

Let \(C_r,F_t\) be those actual factors, and let \(E_r,E_t\)
be arbitrary compatible equal-pair products in their original positions.
Define \(H_{\rm full}=E_rC_rT^L(E_tF_t)\) and
\(I_q(X)=\E[q^{-1/3}|X|^{8/3}]\). Theorem~\ref{thm:odd-two:roots}
gives the complete all-length conclusion
\fitdisplay{
 \E H_{\rm full}^4\le e^{-6d},\qquad
 I_q(T^MH_{\rm full})^3\le e^{-24d/19}.
}{Q192.A.1}
Its proof keeps both finite \(G_1\) boundaries, the full constant
vectors, and the distinct unit and nonunit terminal directions.
The ordinary fourth moment is followed by the paid unanchored outlet
with the original denominator \(q\); no reference arm is added.

Section~\ref{sec:five-chain} assembles these local interfaces with
the appropriate central estimate. The common theorem concerns the
actual local half-tree; its use leaves every root direction, original edge and decoration in place.
At \(d=0\), original leaf integration gives equality; the kernel
itself need not be zero.
